\documentclass[Total.tex]{subfiles}
\begin{document}
	\chapter{Metric Spaces}
	\subsubsection{Tangent Cone as Limit of Cone}
		\begin{Def}
			\begin{itemize}
				\item We could do \textbf{homothety} in the affine space. 
				
				Thus,
				\begin{gather*}
					H_{p,\lambda}：\RR^n\rightarrow \RR^n\qquad x\mapsto \lambda(x-p)
				\end{gather*}
			\end{itemize}
		\end{Def}
		
	\subsubsection{Aymptotic Cone}
	\section{Lipschitz Distance}
		
	\section{Hausdorff-Gromov Distance}
		\subsection{Hausdorff Distance}
			\begin{Def}
				Let:$A,B$ be two metric space. Define:
				\begin{gather*}
				d_H(A,B):=\inf\{\epsilon>0|A\subseteq B_\epsilon(B),B\subseteq B_\epsilon (A)\}
			\end{gather*}
			
			where 
			\begin{gather*}
				B(x,X)=\{x\in X|d(x,A)<\epsilon\}
			\end{gather*}
			\end{Def}
			
			Especially:
			\begin{itemize}
				\item $d_H(A,B)=\infty$ is possible: 
				\begin{gather*}
					A=\{(0,1)\}\qquad B=\{(x,0)|x\in \RR\}
				\end{gather*}
				
				If $\epsilon$ is finite, it is impossible that $B\subseteq B_\epsilon(A)$.
				\item $d_H(A,B),A\cap B$ is possible:
				\begin{gather*}
					Q=\QQ\qquad B=\RR-\QQ
				\end{gather*}
			\end{itemize}
			
			\begin{lemma}
				Suppose $A,B\subseteq X$, and
				\begin{gather*}
					d_A=diam(A)\qquad d_B=diam(B)\\
					d(A,B):=\inf\{d(a,b)|a\in A,b\in B\}
				\end{gather*}
				
				Then,
				\begin{gather*}
					A\subset B_{d_A+d_B+d(A,B)+1}(B)\qquad 
					B\subset B_{d_A+d_B+d(A,B)+1}(A)
				\end{gather*}
			\end{lemma}
			\begin{proof}
				We will expand one of them: By triangle inequality
				\begin{gather*}
					A\subseteq B_{d_A+d(A,B)}(B)\qquad B\subseteq A_{d_A+d(A,B)}(B)
				\end{gather*}
				
				Then, we obtain the relation. Notice: A extreme condition is $d(B)=0$ for pointed set.
			\end{proof}
			
			\begin{corollary} If $A,B$ are bounded(or compact), then
				\begin{gather*}
					d_H(A,B)<\infty
				\end{gather*}
			\end{corollary}
			\begin{proof}
				Omitted.
			\end{proof}
			
			\begin{proposition}
				$d_H$ is a metric on the set of compact subsets.
			\end{proposition}
			\begin{proof}
				\begin{itemize}
					\item Bounded: Thus the value of $d_H$ is in $[0,\infty)$;
					\item Positivity: $d_H(A,A)=0$.
					
					If $d_H(A,B)=0$, then $A\subseteq B,B\subseteq A$;
					\item Symmetry: $d_H(A,B)=d_H(B,A)$
					\item Triangle inequality: $d_H(A,B)=\epsilon_1,d_H(B,C)=\epsilon_2,d_H(A,C)=\epsilon$, then
					\begin{gather*}
						\epsilon_1+
					\end{gather*}
				\end{itemize}
			\end{proof}
			
			We observe: Compactness is a relative weaker condition. We could generalize the theory to 
			\begin{proposition}
				Let $(X,d)$ be a metric space, then, $d_H$ is a metric on 
				\begin{gather*}
					\mathcal{CB}(X)=\{A\subseteq X|A\not=\emptyset,\mbox{ closed and bounded}\}
				\end{gather*}
			\end{proposition}
			\begin{proof}
				Similarly
			\end{proof}
			
			We will take the notion of $d_H^d$ for the metric induced by $d$.
			
			
		\subsection{Gromov-Hausdorff Distance}
		
		\begin{Def}
			\begin{itemize}
				\item $f:(X_1,d_1)\rightarrow (Z,d)$ is called an \textbf{isometric} embedding, if 
				\begin{gather*}
					\bar{f}:(X_1,d_1)\rightarrow (f(X_1),d_2)
				\end{gather*}
				
				is an isometry.
				\item Given two metric space $(X_1,d_1);(X_2,d_2)$, then
				\begin{gather*}
					d_{GH}((X_1,d_1),(X_2,d_2))=
				\end{gather*}
				\item Given two metric space $(X_1,d_1),(X_2,d_2)$, then, the metric space $(Z,d)$ induced by them is 
				\begin{gather*}
					..
				\end{gather*}
			\end{itemize}
		\end{Def}
		
		\begin{proposition}
			Let $(X_1,d_1),(X_2,d_2)$ be metric spaces, then
			\begin{gather*}
				d_{GH}((X_1,d_1),(X_2,d_2))=\inf\{d_H^d(X_1,X_2)\}
			\end{gather*}
		\end{proposition}
		\begin{proof}
			\begin{itemize}
				\item $\leq$:
				
				
				\item $\geq$:
			\end{itemize}
		\end{proof}
		
		\subsection{$\epsilon$-Approxiamation}
			\begin{Def}
				Let:$(X,d_X),(Y,d_Y)$ be two metric spaces. 
				
				A map $f:X\rightarrow Y$ is called $\epsilon$-\textbf{approxiamation} if:
				\begin{itemize}
					\item 
					\item
				\end{itemize}
			\end{Def}
			
			\begin{theorem}
				If $d_{GH}((X_1,d_1),(X_2,d_2))<\epsilon$, then we obtain that, there exists a pair $f,g$
				\begin{gather*}
					Inhalt...
				\end{gather*}
				
				which is an $\epsilon$-approxiamation. i.e. 
				\begin{itemize}
					\item 
					\item 
					\item 
					\item 
				\end{itemize}
			\end{theorem}
			\textbf{Remark} Such given $f,g$ could be not continuous: An example is
			\begin{gather*}
				x_1=S^1,x_2
			\end{gather*}
			\begin{proof}
				Inhalt...
			\end{proof}
			
			Conversely:
			\begin{proposition}
				If the four conditions above are satisfied, then
				\begin{gather*}
					Inhalt...
				\end{gather*}
			\end{proposition}
			
			\subsection{Isometry, Class of Isometric Spaces, and its Gromov-Hausdorff Metric}
				\begin{Def}
					\begin{itemize}
						\item 
						\item 
					\end{itemize}
				\end{Def}
				
				\begin{theorem}
					(Gromov-Hausdorff Distance is a Metric,David A. Edwards,1975)
					
					 The Gromov-Hausdorff distance defines a metric.
				\end{theorem}
				\begin{proof}
					Inhalt...
				\end{proof}
			
			\subsection{General-Topologic Properties}
				\subsubsection{Seperability}
				\subsubsection{Compeleteness}
		\chapter{Topics in Riemannian Geometry}
\end{document}